% part: intuitionistic-logic
% chapter: soundness-completeness
% section: canonical-model

\documentclass[../../../include/open-logic-section]{subfiles}

\begin{document}

\olfileid[id]{int}{sc}{mod}

\olsection{Model Kanonis}

Dunia-dunia dalam model kita akan berupa barisan berhingga~$\sigma$ dari
bilangan asli, yakni $\sigma\in\Nat^*$. Perhatikan bahwa $\Nat^*$
didefinisikan secara induktif oleh:
\begin{enumerate}
\item $\emptyseq \in \Nat^*$.
\item Jika $\sigma\in\Nat^*$ dan $n\in\Nat$, maka
  $\sigma.n\in\Nat^*$ (dengan $\sigma.n$ berarti
  $\sigma\concat\tuple{n}$, sedangkan $\sigma\concat\sigma'$ merupakan
  konkatenasi barisan $\sigma$ dan~$\sigma'$).
\item Tidak ada objek lain yang termasuk dalam $\Nat^*$.
\end{enumerate}
Jadi, kita dapat menggunakan $\Nat^*$ untuk memberikan definisi-definisi
induktif.

Misalkan $\tuple{!B_0,!C_0}$, $\tuple{!B_1,!C_1}$, \dots, merupakan
enumerasi semua pasangan !!{formula}s. Dengan diberikan suatu himpunan
!!{formula}s~$\Delta$, definisikan $\Delta(\sigma)$ secara induktif sebagai
berikut:
\begin{enumerate}
\item $\Delta(\emptyseq)=\Delta$.
\item $\Delta(\sigma.n)={}$
  \[
  \begin{cases}
    (\Delta(\sigma)\cup\{!B_n\})^* &
    \text{jika $\Delta(\sigma)\cup\{!B_n\}\Proves/ !C_n$} \\
    \Delta(\sigma) & \text{jika tidak}.
  \end{cases}
  \]
\end{enumerate}
Di sini, $(\Delta(\sigma)\cup\{!B_n\})^*$ berarti himpunan prima
!!{formula}s yang keberadaannya dijamin oleh
\olref[lin]{lem:lindenbaum}, diterapkan pada himpunan
$\Delta(\sigma)\cup\{!B_n\}$ dan !!{formula}~$!C_n$. Perhatikan bahwa,
menurut definisi ini, jika
$\Delta(\sigma)\cup\{!B_n\}\Proves/ !C_n$, maka
$\Delta(\sigma.n)\Proves !B_n$ dan
$\Delta(\sigma.n)\Proves/ !C_n$. Perhatikan pula bahwa
$\Delta(\sigma)\subseteq\Delta(\sigma.n)$ untuk setiap~$n$. Jika $\Delta$
prima, maka $\Delta(\sigma)$ prima untuk setiap~$\sigma$.

\begin{defn}\ollabel{defn:canonical-model}
  Andaikan $\Delta$ prima. Maka \emph{model kanonis}
  $\mModel{M(\Delta)}$ bagi $\Delta$ didefinisikan oleh:
  \begin{enumerate}
  \item $W=\Nat^*$, yakni himpunan barisan berhingga bilangan asli.
  \item $R$ merupakan tatanan parsial yang memenuhi
    $R\sigma\sigma'$ jika dan hanya jika $\sigma$ merupakan segmen awal
    dari~$\sigma'$ (yakni, $\sigma'=\sigma\concat\sigma''$ untuk suatu
    barisan~$\sigma''$).
  \item $V(p)=\Setabs{\sigma}{p\in\Delta(\sigma)}$.
  \end{enumerate}
\end{defn}

Mudah diperiksa bahwa $R$ memang merupakan tatanan parsial. Syarat
monotonisitas pada~$V$ juga terpenuhi. Karena
$\Delta(\sigma)\subseteq\Delta(\sigma.n)$, dengan induksi pada panjang
ekor~$\sigma''$ kita memperoleh
$\Delta(\sigma)\subseteq\Delta(\sigma')$ setiap kali
$R\sigma\sigma'$.

\end{document}
